% Источники: raw/IMG_7625.jpeg -- raw/IMG_7628.jpeg.
% Компиляция: tectonic 06.tex
\ifdefined\lessonpreambleloaded
  \expandafter\endinput
\fi
\let\lessonpreambleloaded\empty
\documentclass[a4paper,11pt]{article}
\usepackage{fontspec}
\setmainfont{Times New Roman}
\usepackage[russian]{babel}
\usepackage{amsmath,amssymb}
\usepackage[margin=22mm]{geometry}
\usepackage{indentfirst}
\usepackage{needspace}
\usepackage{tikz}
\usetikzlibrary{patterns}
\usepackage[hidelinks]{hyperref}
\usepackage{fancyhdr}
\pagestyle{fancy}
\fancyhf{}
\renewcommand{\headrulewidth}{0pt}
\fancyfoot[C]{\footnotesize\url{https://github.com/bmstu-software-engineering/cs-s5-tv-26y}}
\fancyfoot[R]{\thepage}
\setlength{\emergencystretch}{2em}
\newcommand{\sourcepage}[1]{\par\Needspace{6\baselineskip}\medskip\noindent{\footnotesize\itshape Источник: \detokenize{#1}.}\par\nopagebreak\smallskip}
\newcommand{\editorial}[1]{\par\smallskip\noindent{\small\itshape Примечание к переносу. #1}\par\smallskip}
\newcommand{\lessontitle}[3]{\begin{center}{\large\bfseries Теория вероятностей. #1}\\[4pt]{\large #2}\\[4pt]#3\end{center}}
\newcommand{\parttitle}[1]{\par\Needspace{4\baselineskip}\medskip\noindent\textbf{#1}\par\nopagebreak\smallskip}
\tikzset{element/.style={draw,rectangle,minimum width=7mm,minimum height=5mm,fill=white},every picture/.style={line width=.5pt}}

\begin{document}
\lessontitle{Семинар №\,6}{Комбинаторные задачи}{10.10.26}

\sourcepage{IMG_7625.jpeg}
\section*{Детали: стандартные и нестандартные}
Пронумеруем детали: $1$--$4$ --- нестандартные, $5$--$50$ --- стандартные.
Выбираются $10$ деталей. Порядок не важен, повторений нет.
\editorial{Полное словесное условие на фотографии не приведено;
числа и события ниже взяты из записанного решения.}
Исход --- набор $\{x_1,\ldots,x_{10}\}$ различных номеров из $\{1,\ldots,50\}$:
\[
 \Omega=\bigl\{\{x_1,\ldots,x_{10}\}:x_i\in\{1,\ldots,50\},\
 x_i\ne x_j\text{ при }i\ne j\bigr\}.
\]
\[
 N=|\Omega|=C_{50}^{10}=\frac{50!}{10!40!}.
\]
\parttitle{1. Ровно одна нестандартная деталь}
Пусть $A$ --- это событие. Так как порядок не важен, для описания набора
можно обозначить нестандартную деталь через $x_1$:
\[
 x_1\in\{1,2,3,4\},\qquad x_i\in\{5,\ldots,50\}\quad(i=2,\ldots,10),
\]
причём стандартные детали попарно различны. По основной формуле комбинаторики
\[
 N_A=C_4^1C_{46}^9
 =\frac{4!}{1!3!}\frac{46!}{9!37!}.
\]
\editorial{В одной строке исходника диапазон стандартных деталей начинается с $4$;
здесь исправлено на $5$ согласно нумерации в начале страницы.
Длинные сокращения факториалов зачёркнуты, окончательная вероятность не записана.}

\sourcepage{IMG_7626.jpeg}
\parttitle{2. Хотя бы одна нестандартная деталь}
$B$ --- событие «есть хотя бы одна нестандартная деталь»;
$\overline B$ --- «все детали стандартные».
\[
 P(B)=1-P(\overline B).
\]
Для $\overline B$ выбираем $10$ различных деталей из $46$ стандартных:
\[
 N_{\overline B}=C_{46}^{10}=\frac{46!}{10!36!},
\]
\[
 P(\overline B)=\frac{C_{46}^{10}}{C_{50}^{10}}
 =\frac{46!}{10!36!}\frac{10!40!}{50!}.
\]

\section*{Друзья за столом}
Группа из восьми человек занимает стол. Какова вероятность того,
что два определённых человека окажутся сидящими рядом?
Пронумеруем места от $1$ до $8$. Исход --- перестановка номеров людей:
\[
 \Omega=\{(x_1,\ldots,x_8):x_i\in\{1,\ldots,8\},\
 x_i\ne x_j\text{ при }i\ne j\}.
\]
Порядок важен, повторений нет: это размещения без повторений из $8$ по $8$.

\sourcepage{IMG_7627.jpeg}
\[
 N=A_8^8=\frac{8!}{0!}=8!.
\]
Для двух определённых людей возможны два взаимных порядка; оставшиеся
шесть человек занимают остальные места $6!$ способами.
В первой записи на странице получено
\[
 N_B=8\cdot2\cdot6!,\qquad P(B)=\frac{8\cdot2\cdot6!}{8!}=\frac27.
\]
Ниже отдельно рассматриваются крайние места $1$ и $8$ и внутренние места $2$--$7$.
Для прямого ряда мест записан другой подсчёт:
\[
 P(B)=\frac{(6\cdot2+2\cdot1)\,6!}{8!}=\frac14.
\]
\editorial{В исходнике есть рисунки круглого и прямоугольного стола.
Первый подсчёт учитывает соседство крайних мест друг с другом и соответствует
круговому расположению. Второй --- прямому ряду.
Эти модели нельзя смешивать; исходное словесное условие форму рассадки не уточняет.
Последняя дробь $1/4$ --- упрощение записанной формулы.}

\sourcepage{IMG_7628.jpeg}
\section*{Ботинки}
В шкафу $10$ пар ботинок. Выбирают $4$ ботинка. Найти вероятность того,
что среди них не будет пары.

Порядок не важен, повторений нет. Ботинок обозначим $(n,m)$,
где $n\in\{1,\ldots,10\}$ --- номер пары,
$m\in\{\text{Л},\text{П}\}$ --- левый или правый.
Исход --- набор из четырёх различных ботинок:
\[
 \Omega=\{\{x_1,x_2,x_3,x_4\}:x_i\in
 \{1,\ldots,10\}\times\{\text{Л},\text{П}\},\
 x_i\ne x_j\text{ при }i\ne j\}.
\]
\[
 N=C_{20}^4=\frac{20!}{4!16!}.
\]
Пусть $A$ --- событие «нет ни одной пары». Тогда номера пар у всех четырёх
ботинок различны. Выбираем четыре номера из десяти $C_{10}^4$ способами;
для каждого номера выбираем левый или правый ботинок:
\[
 N_A=C_{10}^4\,2^4,\qquad
 P(A)=\frac{C_{10}^4\,2^4}{C_{20}^4}.
\]
\end{document}
